\documentclass[11pt,a4paper]{article}

\usepackage[margin=2.5cm]{geometry}
\usepackage{amsmath,amssymb}
\usepackage{graphicx}
\usepackage{booktabs}
\usepackage{longtable}
\usepackage{array}
\usepackage{siunitx}
\usepackage{hyperref}
\usepackage{xcolor}
\usepackage{caption}
\usepackage{float}
\usepackage{enumitem}
\usepackage{fancyhdr}
\usepackage{titling}

\hypersetup{
    colorlinks=true,
    linkcolor=blue!50!black,
    urlcolor=blue!50!black,
    citecolor=blue!50!black
}

\pagestyle{fancy}
\fancyhf{}
\fancyhead[L]{Temperature Sensors: Calibration of Pt-100, Diode and Thermocouple}
\fancyhead[R]{\thepage}
\renewcommand{\headrulewidth}{0.4pt}

\setlength{\parindent}{0pt}
\setlength{\parskip}{6pt}

\title{\textbf{Temperature Sensors} \\[4pt] \large Calibration of a Pt-100 Resistance Thermometer, a Semiconductor Diode, and a Thermocouple}
\author{Parth Soni}
\date{\today}

\begin{document}

\begin{titlepage}
\centering
\vspace*{2cm}
{\Huge \textbf{Temperature Sensors}}\\[0.5cm]
{\Large Calibration of Pt-100, Semiconductor Diode, and Thermocouple\\ as Temperature Sensors}\\[2cm]

{\large \textbf{Submitted by}}\\[0.3cm]
{\large Parth Soni}\\[2cm]

{\large \textbf{Experiment:} Temperature Sensors}\\[0.3cm]
{\large \textbf{Date of experiment:} 24 August 2026}\\[0.3cm]

\vfill
{\large \today}
\end{titlepage}

\tableofcontents
\newpage

\section{Aim of the Experiment}
\begin{enumerate}[label=(\roman*)]
    \item To calibrate a Pt-100 platinum resistance thermometer using the boiling point of liquid nitrogen, the melting point of ice, room temperature, and the boiling point of water as fixed/standard temperature references, and to obtain the relation between resistance and temperature for Pt-100.
    \item To use the calibrated Pt-100 as a reference thermometer and calibrate a semiconductor p--n junction diode (1N4148) as a temperature sensor by measuring its forward voltage as a function of temperature, and to estimate the band gap energy of the semiconductor.
    \item To calibrate a Chromel--Alumel thermocouple as a temperature sensor by measuring the thermo-emf developed as a function of temperature.
\end{enumerate}

\section{Theory}

\subsection{General idea of temperature measurement}
Temperature cannot be measured directly like length or mass; it is inferred from some other physical property of a system that changes monotonically and reproducibly with temperature, and which can be measured directly (length of a mercury column, electrical resistance, voltage across a junction, etc.). In this experiment three different physical effects are used as the basis of three different classes of temperature sensors: (i) the temperature dependence of the electrical resistance of a metal (Pt-100), (ii) the temperature dependence of the forward voltage of a semiconductor p--n junction (diode), and (iii) the Seebeck effect in a junction of two dissimilar metals (thermocouple). Since the Pt-100 sensor can be calibrated against well-defined fixed points (boiling point of liquid nitrogen, melting point of ice, boiling point of water, room temperature), it is used here as the \emph{reference thermometer} against which the other two sensors are calibrated.

\subsection{Pt-100 Resistance Thermometer}
The electrical resistivity of a pure metal arises from the scattering of conduction electrons by lattice vibrations (phonons) and by impurities/defects. As temperature increases, the amplitude of lattice vibrations increases, the phonon scattering rate increases, and hence the resistivity -- and therefore the resistance of a wire of fixed geometry -- increases. Over a limited range of temperature this variation is very nearly linear and can be written as
\begin{equation}
    R(T) = R_0\left(1 + \alpha T\right),
    \label{eq:pt100}
\end{equation}
where $R_0$ is the resistance at $T = 0\,^{\circ}\mathrm{C}$, $\alpha$ is the temperature coefficient of resistance of the metal, and $T$ is the temperature in $^{\circ}\mathrm{C}$. For a standard Pt-100 sensor $R_0 = 100\,\Omega$ and $\alpha \approx 3.85\times10^{-3}\,^{\circ}\mathrm{C}^{-1}$. Equation~\eqref{eq:pt100} can equivalently be written as a straight line
\begin{equation}
    R(T) = a + bT, \qquad a = R_0, \quad b = R_0\alpha,
    \label{eq:pt100linear}
\end{equation}
so that a linear regression of measured $R$ against the known standard temperatures $T$ directly yields $R_0$ and $\alpha$ for the particular sensor used. Once $a$ and $b$ are known, Eq.~\eqref{eq:pt100linear} can be inverted to obtain the unknown temperature from a measured resistance:
\begin{equation}
    T(R) = \frac{R - a}{b}.
    \label{eq:pt100inv}
\end{equation}
Platinum is chosen because it is chemically inert (does not oxidise easily), has a highly linear and reproducible $R$--$T$ characteristic over a wide range, and can be obtained in a very pure form, which is essential for reproducibility from one sensor to another.

\subsection{Semiconductor Diode as a Temperature Sensor}
In a p--n junction diode under forward bias, the current--voltage relation is given by the ideal diode equation
\begin{equation}
    I = I_0\left(e^{qV/\eta k T} - 1\right),
    \label{eq:diode1}
\end{equation}
where $T$ is the absolute temperature of the junction, $q$ is the electronic charge, $k$ is Boltzmann's constant, and $\eta$ (the ideality factor) lies between 1 and 2. The saturation current $I_0$ itself depends strongly on temperature and on the band gap energy $E_g$ of the semiconductor:
\begin{equation}
    I_0 = p\,T^{m}\, e^{-E_g/\eta k T},
    \label{eq:I0}
\end{equation}
where $p$ is a constant depending on the junction geometry and doping, and $m$ is a material-dependent constant.

For a forward-biased junction with $V \gg \eta k T/q$, the $-1$ in Eq.~\eqref{eq:diode1} may be dropped:
\begin{equation}
    I \simeq I_0\, e^{qV/\eta k T}.
    \label{eq:diode2}
\end{equation}
Substituting Eq.~\eqref{eq:I0} into Eq.~\eqref{eq:diode2},
\begin{equation}
    I = p\,T^{m}\, e^{-E_g/\eta k T}\, e^{qV/\eta k T}
      = p\,T^{m}\, \exp\!\left[\frac{qV - E_g}{\eta k T}\right].
    \label{eq:diode3}
\end{equation}

\subsubsection*{Derivation of $V(T)$}
Taking the natural logarithm of Eq.~\eqref{eq:diode3},
\begin{equation}
    \ln I = \ln p + m\ln T + \frac{qV - E_g}{\eta k T}.
\end{equation}
Rearranging for $qV$,
\begin{equation}
    qV = E_g + \eta k T\left[\ln I - \ln p - m\ln T\right]
       = E_g - \eta k T\left[m\ln T + \ln p - \ln I\right].
\end{equation}
Defining
\begin{equation}
    C \equiv m\ln T + \ln p - \ln I,
    \label{eq:Cdef}
\end{equation}
we obtain
\begin{equation}
    qV = E_g - C\eta k T
    \qquad\Longrightarrow\qquad
    \boxed{V = \frac{E_g}{q} - \frac{C\eta k}{q}\,T}.
    \label{eq:diodeVT}
\end{equation}
If the diode is operated at constant forward current $I$ (as is done in this experiment), and since the variation of $\ln T$ with $T$ in Eq.~\eqref{eq:Cdef} is very slow (logarithmic) compared to the much stronger linear variation of $V$ with $T$ contained in the second term of Eq.~\eqref{eq:diodeVT}, $C$ may be treated as effectively constant. Equation~\eqref{eq:diodeVT} then predicts that \emph{the forward voltage of a diode operated at constant current varies linearly with the absolute temperature $T$}, with a negative slope; this is the basis for using a diode as a temperature sensor. Writing
\begin{equation}
    V(T) = A - B\,T, \qquad A = \frac{E_g}{q}, \quad B = \frac{C\eta k}{q},
    \label{eq:diodefit}
\end{equation}
a linear fit of the measured $V$--$T$ data gives the slope $B$ and the intercept $A$. Since $C\approx \mathcal{O}(1)$ and $\eta \approx 1$--$2$, the intercept $A$ obtained by extrapolating the straight line to $T=0\,\mathrm{K}$ gives an order-of-magnitude estimate of the band-gap energy of the semiconductor, $E_g \approx qA$ (in eV, $E_g$ is numerically equal to $A$ expressed in volts).

\subsection{Thermocouple}
When two dissimilar metals (here Chromel and Alumel) are joined to form a junction, and this junction is held at a temperature $T$ different from a reference junction (here, formed internally at the measuring instrument, effectively at room temperature $T_{\mathrm{ref}}$), a small electromotive force (thermo-emf) is developed between the free ends. This is known as the \emph{Seebeck effect}, and arises because the density of free electrons and their diffusion characteristics differ between the two metals; when a temperature gradient exists across each conductor, electrons diffuse from the hot end to the cold end at different rates in the two metals, setting up a net contact potential difference at the junction that is not exactly cancelled between the two metals used.

For a single junction thermocouple and over a limited temperature range around the reference point, the relationship between the thermo-emf $\varepsilon$ and the junction temperature $T$ is very nearly linear (to lowest order in a Taylor expansion of the true, weakly non-linear characteristic):
\begin{equation}
    \varepsilon(T) = \varepsilon_0 + S\,T,
    \label{eq:thermocouple}
\end{equation}
where $S = d\varepsilon/dT$ is called the \emph{Seebeck coefficient} (or thermoelectric sensitivity) of the thermocouple pair, and $\varepsilon_0$ is the (small, ideally zero) emf recorded when the sensing junction is at the reference-junction temperature. A linear regression of measured emf against the (independently known, here Pt-100-derived) temperature gives $S$ and $\varepsilon_0$, which together constitute the calibration of the thermocouple.

\section{Apparatus}
Pt-100 platinum resistance sensor; digital multimeter (four-probe resistance measurement); semiconductor diode 1N4148 operated at constant forward current, with a digital voltmeter; Chromel--Alumel thermocouple with a millivoltmeter; liquid nitrogen (LN$_2$) in a thermocol (Dewar) container; ice; a beaker of water with a heater for the boiling-point measurement; thermometer / room-temperature reference.

\section{Procedure}
\begin{enumerate}
    \item The resistance of the Pt-100 sensor was measured at four standard/reference temperatures: room temperature ($24.8\,^{\circ}\mathrm{C}$, measured with a laboratory thermometer), melting point of ice ($0.5\,^{\circ}\mathrm{C}$), boiling point of water at the local atmospheric pressure ($96.1\,^{\circ}\mathrm{C}$), and the boiling point of liquid nitrogen ($-196\,^{\circ}\mathrm{C}$). A linear fit of $R$ vs.\ $T$ then gave the calibration equation for the Pt-100 [Eq.~\eqref{eq:pt100linear}].
    \item The Pt-100, the diode, and the thermocouple junction were then bundled together and immersed in the liquid-nitrogen vapour above the LN$_2$ surface (which has a temperature gradient from $\approx -196\,^{\circ}\mathrm{C}$ near the surface to room temperature higher up), so that all three sensors are always at (nearly) the same temperature at any instant.
    \item The assembly was allowed to warm up slowly from liquid-nitrogen temperature to room temperature. At each instant, the resistance $R$ of the Pt-100, the thermo-emf $\varepsilon$ of the thermocouple, and the forward voltage $V$ of the diode (at fixed forward current) were recorded simultaneously.
    \item The Pt-100 resistance $R$ recorded at each instant was converted into the corresponding temperature $T$ using the inverted calibration equation~\eqref{eq:pt100inv} (and $T[\mathrm{K}] = T[^\circ\mathrm{C}] + 273.15$ where an absolute temperature was needed, i.e.\ for the diode).
    \item The thermo-emf $\varepsilon$ and the diode voltage $V$ were then each plotted against the derived temperature $T$, and linear regression was used to obtain the respective calibration equations [Eqs.~\eqref{eq:diodefit} and \eqref{eq:thermocouple}].
\end{enumerate}

\section{Observations / Data}

\subsection{Pt-100 calibration data}
\begin{table}[H]
\centering
\caption{Pt-100 resistance at four standard reference temperatures.}
\begin{tabular}{lcc}
\toprule
Reference point & $T$ ($^{\circ}$C) & $R$ ($\Omega$) \\
\midrule
Room temperature   & 24.8   & 110.1 \\
Boiling point of water & 96.1  & 136.2 \\
Melting point of ice   & 0.5   & 100.8 \\
Boiling point of LN$_2$ & $-196$ & 20.0 \\
\bottomrule
\end{tabular}
\end{table}

\subsection{Pt-100 resistance, thermocouple emf and diode voltage (simultaneous readings during warm-up)}
\begin{table}[H]
\centering
\caption{Simultaneously recorded Pt-100 resistance $R$, thermocouple emf $\varepsilon$, and diode forward voltage $V$, along with the temperature $T$ derived from $R$ using Eq.~\eqref{eq:pt100inv}.}
\begin{tabular}{cccccc}
\toprule
$R$ ($\Omega$) & $T$ ($^\circ$C) & $T$ (K) & $\varepsilon$ (mV) & $V$ (V) \\
\midrule
20.0 & $-197.65$ & 75.50 & 6.4 & 1.020 \\
29.2 & $-174.70$ & 98.45 & 6.3 & 1.010 \\
30.2 & $-172.21$ & 100.94 & 6.3 & 1.010 \\
35.0 & $-160.24$ & 112.91 & 6.1 & 1.000 \\
40.0 & $-147.77$ & 125.38 & 5.9 & 0.990 \\
45.9 & $-133.06$ & 140.09 & 5.7 & 0.970 \\
50.1 & $-122.58$ & 150.57 & 5.4 & 0.950 \\
55.0 & $-110.36$ & 162.79 & 5.1 & 0.930 \\
60.0 & $-97.89$ & 175.26 & 4.7 & 0.910 \\
65.0 & $-85.42$ & 187.73 & 4.3 & 0.890 \\
70.0 & $-72.95$ & 200.20 & 3.9 & 0.870 \\
75.0 & $-60.48$ & 212.67 & 3.5 & 0.850 \\
80.0 & $-48.02$ & 225.13 & 3.1 & 0.830 \\
85.2 & $-35.05$ & 238.10 & 2.6 & 0.820 \\
90.0 & $-23.08$ & 250.07 & 2.2 & 0.800 \\
96.4 & $-7.12$ & 266.03 & 1.5 & 0.770 \\
100.0 & 1.86 & 275.01 & 1.1 & 0.760 \\
105.0 & 14.33 & 287.48 & 0.6 & 0.740 \\
109.6 & 23.31 & 296.46 & 0.1 & 0.720 \\
\bottomrule
\end{tabular}
\end{table}

\section{Calculations and Analysis}

\subsection{Pt-100 calibration}
A least-squares straight-line fit $R = a + bT$ was carried out on the four points of Table~1 using
\begin{equation}
    b = \frac{\sum (T_i-\bar T)(R_i - \bar R)}{\sum (T_i - \bar T)^2}, \qquad a = \bar R - b\bar T ,
\end{equation}
which gave
\begin{equation}
    \boxed{R(T) = (99.25 \pm 0.85) + (0.4010 \pm 0.0077)\,T \ \ [\Omega,\ ^\circ\mathrm{C}]}
\end{equation}
with correlation coefficient $r = 0.9996$ ($r^2 = 0.9993$). Comparing with Eq.~\eqref{eq:pt100linear}, this gives $R_0 = a = 99.25\pm0.85\,\Omega$ and
\begin{equation}
    \alpha = \frac{b}{a} = (4.04 \pm 0.08)\times 10^{-3}\ ^\circ\mathrm{C}^{-1},
\end{equation}
which is close to the standard reference value $\alpha_{\text{Pt}} \approx 3.85\times10^{-3}\ ^\circ\mathrm{C}^{-1}$, and $R_0 \approx 100\,\Omega$ as expected for a Pt-100 sensor, confirming that the sensor behaves as expected. This calibration line, together with the four data points, is shown in Fig.~\ref{fig:pt100}. The inverse relation used to convert resistance into temperature for the rest of the experiment is
\begin{equation}
    T(R) = \frac{R - 99.25}{0.4010}\ \ [^\circ\mathrm{C}].
\end{equation}

\begin{figure}[H]
\centering
\includegraphics[width=0.72\textwidth]{graphs/pt100_calibration.pdf}
\caption{Calibration of the Pt-100 resistance thermometer: resistance vs.\ temperature at the four standard reference points, with the best-fit straight line.}
\label{fig:pt100}
\end{figure}

\subsection{Semiconductor diode calibration}
Using the temperatures $T$(K) obtained from the Pt-100 (Table~2), the diode forward voltage $V$ was fitted to the linear form of Eq.~\eqref{eq:diodefit}, $V = A - BT$:
\begin{equation}
    \boxed{V(T) = (1.1559 \pm 0.0070)\ \mathrm{V} - (1.433 \pm 0.035)\times10^{-3}\ \mathrm{V/K}\ \times T}
\end{equation}
with $r = -0.9949$ ($r^2 = 0.9899$), i.e.\ the diode voltage decreases essentially linearly with increasing temperature, as predicted by Eq.~\eqref{eq:diodeVT} (Fig.~\ref{fig:diode}). The sensitivity of the diode as a temperature sensor is
\begin{equation}
    \left|\frac{dV}{dT}\right| = 1.433\ \mathrm{mV/K}.
\end{equation}
Extrapolating the fitted line to $T = 0\,\mathrm{K}$ gives the intercept $A = E_g/q$, so that the band-gap energy of the diode material (silicon, for the 1N4148) is estimated as
\begin{equation}
    \boxed{E_g \approx qA = 1.16 \pm 0.01\ \mathrm{eV}},
\end{equation}
which compares well with the accepted band gap of silicon at low temperature, $E_g(\mathrm{Si}) \approx 1.12\ \mathrm{eV}$ (the small excess is expected, since Eq.~\eqref{eq:diodeVT} neglects the weak $\ln T$ dependence of $C$, and $E_g$ itself decreases slightly with increasing $T$ in real silicon, which is not captured by a straight-line extrapolation from the $75$--$300$ K range used here).

\begin{figure}[H]
\centering
\includegraphics[width=0.72\textwidth]{graphs/diode_calibration.pdf}
\caption{Forward voltage of the 1N4148 diode (at constant forward current) as a function of absolute temperature, with best-fit straight line. The intercept gives an estimate of the band gap energy.}
\label{fig:diode}
\end{figure}

\subsection{Thermocouple calibration}
The thermo-emf $\varepsilon$ was fitted to the linear form of Eq.~\eqref{eq:thermocouple}, $\varepsilon = \varepsilon_0 + ST$, using the Pt-100-derived Celsius temperatures:
\begin{equation}
    \boxed{\varepsilon(T) = (1.419 \pm 0.142)\ \mathrm{mV} - (0.0297 \pm 0.0013)\ \mathrm{mV}/^\circ\mathrm{C} \times T}
\end{equation}
with $r = -0.984$ ($r^2 = 0.968$). The magnitude of the Seebeck coefficient obtained is
\begin{equation}
    |S| = \left|\frac{d\varepsilon}{dT}\right| \approx 29.7\pm1.3\ \mu\mathrm{V}/^{\circ}\mathrm{C},
\end{equation}
which is of the same order as the nominal Seebeck coefficient quoted for Chromel--Alumel (Type K) thermocouples ($\approx 40\ \mu\mathrm{V}/^\circ\mathrm{C}$ near room temperature); the somewhat lower value obtained here, and the slightly poorer linearity ($r^2=0.968$, visibly non-zero curvature at the two ends of the range in Fig.~\ref{fig:thermocouple}), reflects the fact that the true Seebeck coefficient of a thermocouple is itself a (weak) function of temperature, so that a single straight line is only an approximation over the very wide range ($\approx-198\,^\circ$C to $+23\,^\circ$C) spanned in this experiment. The emf extrapolates to zero close to room temperature, consistent with the reference (cold) junction of the digital millivoltmeter being at room temperature.

\begin{figure}[H]
\centering
\includegraphics[width=0.72\textwidth]{graphs/thermocouple_calibration.pdf}
\caption{Thermo-emf of the Chromel--Alumel thermocouple as a function of temperature, with best-fit straight line.}
\label{fig:thermocouple}
\end{figure}

\section{Error Analysis}

\subsection{Method}
For each calibration $y = a + bx$, the least-squares estimates of slope and intercept and their standard errors were computed from
\begin{equation}
    b = \frac{S_{xy}}{S_{xx}}, \quad a = \bar y - b\bar x, \quad
    s = \sqrt{\frac{\sum(y_i - a - bx_i)^2}{n-2}},
\end{equation}
\begin{equation}
    \delta b = \frac{s}{\sqrt{S_{xx}}}, \qquad
    \delta a = s\sqrt{\frac{\sum x_i^2}{n\, S_{xx}}},
\end{equation}
where $S_{xx} = \sum (x_i - \bar x)^2$, $S_{xy}=\sum(x_i-\bar x)(y_i - \bar y)$, and $n$ is the number of data points. The goodness of the linear fit is quantified by the Pearson correlation coefficient $r = S_{xy}/\sqrt{S_{xx}S_{yy}}$.

\subsection{Summary of uncertainties}
\begin{table}[H]
\centering
\caption{Summary of least-squares fit parameters and their standard errors for the three sensors.}
\begin{tabular}{lcccc}
\toprule
Sensor & Fit & Slope & Intercept & $r^2$ \\
\midrule
Pt-100 & $R=a+bT$ ($\Omega$, $^\circ$C) & $0.4010\pm0.0077$ & $99.25\pm0.85$ & 0.9993 \\
Diode  & $V=A+BT$ (V, K)  & $(-1.433\pm0.035)\times10^{-3}$ & $1.1559\pm0.0070$ & 0.9899 \\
Thermocouple & $\varepsilon=\varepsilon_0+ST$ (mV, $^\circ$C) & $-0.0297\pm0.0013$ & $1.419\pm0.142$ & 0.9682 \\
\bottomrule
\end{tabular}
\end{table}

\subsection{Sources of error}
\begin{itemize}
    \item \textbf{Pt-100 calibration itself rests on only four points}, one of which (the LN$_2$ boiling point) is obtained from a vapour region above the liquid whose temperature can drift with the LN$_2$ level and with how deep the sensor is immersed; this is the largest single source of uncertainty in $a$ and $b$ (note the largest residual, $-1.6\,\Omega$, occurs at $T=96.1\,^\circ$C, Fig.~\ref{fig:residuals}).
    \item \textbf{Non-simultaneity / thermal lag}: the Pt-100, diode, and thermocouple, although bundled together, do not have identical thermal mass and contact area, so during the warm-up run there can be a small lag between the temperature actually seen by each sensor at a given instant; this contributes scatter to the diode and thermocouple fits over and above pure instrumental error.
    \item \textbf{Instrumental resolution}: the multimeter/millivoltmeter readings are limited by the last significant digit of the display (e.g.\ $\pm0.1\,\mathrm{mV}$ for the thermocouple, $\pm0.001$--$0.01\,\mathrm{V}$ for the diode, $\pm0.1\,\Omega$ for the Pt-100), which sets a floor on the achievable precision.
    \item \textbf{Non-linearity over a wide temperature range}: both the true $R$--$T$ characteristic of platinum and the Seebeck coefficient of the thermocouple are only \emph{approximately} linear; deviations from linearity are most visible at the extremes of the range (very low or near room temperature), and show up as the systematic curvature seen in Fig.~\ref{fig:thermocouple} and as the comparatively lower $r^2$ for the thermocouple fit.
    \item \textbf{Contact/lead resistance} in the two/four-wire measurement of the Pt-100 resistance, and thermal emf's generated at the connecting leads of the thermocouple (parasitic junctions), are systematic errors not accounted for in the statistical error analysis above.
\end{itemize}

\begin{figure}[H]
\centering
\includegraphics[width=0.65\textwidth]{graphs/pt100_residuals.pdf}
\caption{Residuals ($R_{\text{observed}} - R_{\text{fit}}$) of the Pt-100 calibration fit at the four reference temperatures.}
\label{fig:residuals}
\end{figure}

\section{Results and Conclusion}
\begin{enumerate}[label=(\arabic*)]
    \item The Pt-100 sensor was successfully calibrated: $R(T) = (99.25\pm0.85) + (0.4010\pm0.0077)\,T\ [\Omega,\,^\circ\text{C}]$, giving $R_0 \approx 99.3\,\Omega$ and $\alpha \approx 4.04\times10^{-3}\ ^\circ\mathrm{C}^{-1}$, consistent with the standard Pt-100 characteristic ($R_0=100\,\Omega$, $\alpha \approx 3.85\times10^{-3}\ ^\circ\mathrm{C}^{-1}$).
    \item Using this as a reference thermometer, the 1N4148 diode was calibrated over the range $\sim$75--296~K: $V(T) = (1.1559\pm0.0070) - (1.433\pm0.035)\times10^{-3}\,T$ [V, K], a highly linear relation ($r^2=0.99$), from which the semiconductor band gap was estimated as $E_g \approx 1.16\,\mathrm{eV}$, in good agreement with the known band gap of silicon ($\approx1.12$~eV).
    \item The Chromel--Alumel thermocouple was calibrated over the range $\sim{-198}$ to $+23\,^\circ$C: $\varepsilon(T) = (1.419\pm0.142) - (0.0297\pm0.0013)\,T$ [mV, $^\circ$C], with somewhat greater deviation from linearity ($r^2=0.97$) than the other two sensors, reflecting the intrinsically weaker temperature-dependence of the Seebeck coefficient of Chromel--Alumel over such a wide span.
    \item The Pt-100 has the largest fractional/absolute change of its measured quantity per degree, followed by the thermocouple, with the diode being intermediate; all three sensors, however, gave temperature calibrations consistent with expected physics, validating the experimental method.
\end{enumerate}
\centering
{\large \textbf{-End of the report-}}\\[0.3cm]




\end{document}